\subsection{多重模族的张量积}

设 \((\mathbf{G}_{\lambda})_{\lambda \in L}\) 是一族 \(\mathbf{Z}\) -模；一个映射 \(u\) （从集合 \(\mathbf{G} = \prod_{\lambda \in L} \mathbf{G}_{\lambda}\) 到一个 \(\mathbf{Z}\) -模）称为多重加性的（或 \(\mathbf{Z}\) -多重线性的），如果 \((x_{\lambda}) \mapsto u((x_{\lambda}))\) 关于每个变量 \(x_{\lambda}\) 都是加性的；确切地说，这意味着，对于所有 \(\mu \in L\) 以及每个元素 \((a_{\lambda}) \in \prod_{\lambda \neq \mu} \mathbf{G}_{\lambda}\) ，把 \(\mathbf{G}\) 与 \(\mathbf{G}_{\mu} \times \prod_{\lambda \neq \mu} \mathbf{G}_{\lambda}\) 典范地等同起来，

\[
u (x _ {\mu} + y _ {\mu}, (a _ {\lambda})) = u (x _ {\mu}, (a _ {\lambda})) + u (y _ {\mu}, (a _ {\lambda})) \quad \text { for } x _ {\mu}, y _ {\mu} \text { in } G _ {\mu}.\tag{24}
\]

这特别意味着 \(u((x_{\lambda})) = 0\) （若 \(x_{\lambda}\) 中有一个为零）。

我们还考虑如下的泛映射问题，其中 \(\Sigma\) 是 Z-模结构的种类，而 \(\alpha\) -映射是将 G 映入一个 Z-模的多重加性映射。其解仍可通过考虑 Z-模 \(C = Z^{(G)}\) （即 G 的元素以 Z 中系数作形式线性组合所构成的模）以及 C 中由形如

\[
\left(x _ {\mu} + y _ {\mu}, \left(z _ {\lambda}\right) _ {\lambda \neq \mu}\right) - \left(x _ {\mu}, \left(z _ {\lambda}\right) _ {\lambda \neq \mu}\right) - \left(y _ {\mu}, \left(z _ {\lambda}\right) _ {\lambda \neq \mu}\right)
\]

（其中 \(\mu \in \mathbf{L}\) , \(x_{\mu} \in \mathrm{G}_{\mu}\) , \(y_{\mu} \in \mathrm{G}_{\mu}\) 和 \(z_{\lambda} \in \mathrm{G}_{\lambda}\) ( \(\lambda \neq \mu\) ) 都是任意的）的元素生成的子 Z-模 D 来获得。商 \(\mathbf{Z}\) -模 C/D 称为 \(\mathbf{Z}\) -模族 \((\mathrm{G}_{\lambda})_{\lambda \in L}\) 的（在 \(\mathbf{Z}\) 上的）张量积，并记作 \(\bigotimes_{\lambda \in L} \mathrm{G}_{\lambda}\) ；对每一个元素 \((x_{\lambda})_{\lambda \in L}\) （属于 G 且为 C 的典范基的一个元素）， \(\bigotimes_{\lambda \in L} x_{\lambda}\) 表示该元素在 C/D 中的典范像。由上述定义可知，映射 \(\phi: (x_{\lambda}) \mapsto \bigotimes_{\lambda \in L} x_{\lambda}\) （从 G 到 \(\bigotimes_{\lambda \in L} G_{\lambda}\) ）是 Z-多重线性的，并且对于每一个将 G 映入一个 Z-模 H 的 Z-多重线性映射 \(f\) ，存在且仅存在一个 Z-线性映射 \(g: \bigotimes_{\lambda \in L} G_{\lambda} \to H\) ，使得 \(f = g \circ \phi\) ；有序对 \(\left( \bigotimes_{\lambda \in L} G_{\lambda}, \phi \right)\) 从而解决了所讨论的泛映射问题。

设 \((\mathbf{G}_{\lambda}^{\prime})_{\lambda \in L}\) 为 \(\mathbf{Z}\) -模的另一个族，并且对于所有 \(\lambda \in L\) ，设 \(v_{\lambda}: \mathbf{G}_{\lambda} \to \mathbf{G}_{\lambda}^{\prime}\) 是 \(\mathbf{Z}\) -线性映射（换言之，即交换群的同态）。则映射

\[
(x _ {\lambda}) \mapsto \bigotimes_ {\lambda \in L} v _ {\lambda} (x _ {\lambda})
\]

（从 \(\mathbf{G}\) 到 \(\bigotimes_{\lambda \in L} \mathbf{G}_{\lambda}'\) ）是 \(\mathbf{Z}\) -多重线性的，因此典范地定义了一个 \(\mathbf{Z}\) -线性映射（从 \(\bigotimes_{\lambda \in L} \mathbf{G}_{\lambda}\) 到 \(\bigotimes_{\lambda \in L} \mathbf{G}_{\lambda}'\) ），记作 \(\bigotimes_{\lambda \in L} v_{\lambda}\) ，且使得

\[
\left(\bigotimes_ {\lambda \in L} v _ {\lambda}\right) \left(\bigotimes_ {\lambda \in L} x _ {\lambda}\right) = \bigotimes_ {\lambda \in L} v _ {\lambda} (x _ {\lambda}).\tag{25}
\]

特别地，我们考虑，对于某个 \(\mu \in \mathbf{L}\) ，一个自同态 \(\theta\) （ \(\mathbf{G}_{\mu}\) 的）；我们用 \(\tilde{\theta}\) 表示 \(\bigotimes_{\lambda \in \mathbf{L}} \mathbf{G}_{\lambda}\) 的等于 \(\bigotimes_{\lambda \in \mathbf{L}} v_{\lambda}\) 的自同态（其中 \(v_{\mu} = \theta\) ， \(v_{\lambda} = 1_{\mathbf{G}_{\lambda}}\) （对 \(\lambda \neq \mu\) ））。

然后我们假设给定了一个集合 \(\Omega\) 、一个映射

\[
c: \omega \mapsto (\rho (\omega), \sigma (\omega))
\]

（从 \(\Omega\) 到 \(\mathbf{L} \times \mathbf{L}\) ）并且，对于所有 \(\omega \in \Omega\) ，给定一个自同态 \(p_{\omega}\) （ \(\mathrm{G}_{\rho(\omega)}\) 的）以及一个自同态 \(q_{\omega}\) （ \(\mathrm{G}_{\sigma(\omega)}\) 的）；对应于它们有两个自同态 \(\tilde{p}_{\omega}\) 和 \(\tilde{q}_{\omega}\) （属于 \(\mathrm{P} = \bigotimes_{\lambda \in \mathbf{L}} \mathrm{G}_{\lambda}\) ）。设 \(\mathbf{R}\) 为子 \(\mathbf{Z}\) -模（ \(\mathbf{P}\) 中的），由自同态 \(\tilde{p}_{\omega} - \tilde{q}_{\omega}\) （当 \(\omega\) 遍历 \(\Omega\) 时）的像的并集所生成。商 \(\mathbf{Z}\) -模 \(\mathrm{P/R}\) 被称为族 \((\mathrm{G}_{\lambda})_{\lambda \in \mathbf{L}}\) 相对于 \(c, p, q\) 的张量积，并且记作 \(\bigotimes_{(c, p, q)} \mathrm{G}_{\lambda}\) ；将典范同态 \(\mathrm{P} \to \mathrm{P/R}\) 与上面定义的映射 \(\phi: \mathrm{G} \to \bigotimes_{\lambda \in \mathbf{L}} \mathrm{G}_{\lambda}\) 复合，我们得到一个 \(\mathbf{Z}\) -多重线性映射 \(\phi_{(c, p, q)}: \mathrm{G} \to \bigotimes_{(c, p, q)} \mathrm{G}_{\lambda}\) ，并记 \(\phi_{(c, p, q)}((x_{\lambda})) = \bigotimes_{(c, p, q)} x_{\lambda}\) 或简记为 \(\bigotimes_{(c)} x_{\lambda}\) 。由 \(\bigotimes_{(c, p, q)} x_{\lambda}\) 和 \(\phi_{(c, p, q)}\) 组成的有序对解决了以下泛映射问题：设 \(\tilde{p}_{\omega}\) （相应地 \(\tilde{q}_{\omega}\) ）表示从

\[
\left(x _ {\rho (\omega)}, \left(x _ {\lambda}\right) _ {\lambda \neq \rho (\omega)}\right) \mapsto \left(p _ {\omega} \left(x _ {\rho (\omega)}\right), \left(x _ {\lambda}\right) _ {\lambda \neq \rho (\omega)}\right)
\]

\[
\left(\text { resp. } \left(x _ {\sigma (\omega)}, \left(x _ {\lambda}\right) _ {\lambda \neq \sigma (\omega)}\right) \mapsto \left(q _ {\omega} (x _ {\sigma (\omega)}), \left(x _ {\lambda}\right) _ {\lambda \neq \sigma (\omega)}\right)\right)
\]

到自身的映射。那么 \(\Sigma\) 取为 Z-模结构的种类，而 \(\alpha\) -映射取为 \(\mathbf{Z}\) -多重线性映射 \(u\) （从 \(\mathbf{G}\) 到一个 \(\mathbf{Z}\) -模），它们进一步满足条件

\[
u \circ \tilde {p} _ {\omega} = u \circ \tilde {q} _ {\omega}\tag{26}
\]

（对所有 \(\omega \in \Omega\) ）。由上述定义，证明是显然的。

这一构造特别恢复了第 1 号中描述的 \(\mathbf{E} \otimes_{\mathbf{A}} \mathbf{F}\) 的构造：在这种情况下，我们取 \(\mathbf{L} = \{1, 2\}\) , \(\mathbf{G}_1 = \mathbf{E}\) , \(\mathbf{G}_2 = \mathbf{F}\) , \(\Omega = \mathbf{A}\) ；此外，对于所有 \(\omega \in \mathbf{A}\) ，我们必须有 \(\rho(\omega) = 1\) , \(\sigma(\omega) = 2\) , \(p_\omega\) 是自同态 \(x \mapsto x\omega\) （ \(\mathbf{Z}\) -模 \(\mathbf{E}\) 的），而 \(q_\omega\) 是自同态 \(y \mapsto \omega y\) （ \(\mathbf{Z}\) -模 \(\mathbf{F}\) 的）。

设 \((\mathbf{G}_{\lambda}^{\prime})_{\lambda \in L}\) 是 Z-模的第二个族；保持映射 c 不变，假设对所有 \(\omega \in \Omega\) 给定了一个自同态 \(p_{\omega}^{\prime}\) （ \(\mathbf{G}_{\rho(\omega)}^{\prime}\) 的）以及一个自同态 \(q_{\omega}^{\prime}\) （ \(\mathbf{G}_{\sigma(\omega)}^{\prime}\) 的）。然后对所有 \(\lambda \in L\) ，令 \(v_{\lambda}: G_{\lambda} \to G_{\lambda}^{\prime}\) 是一个 Z-线性映射，使得对所有 \(\omega \in \Omega\) ,

\[
v _ {\rho (\omega)} \circ p _ {\omega} = p _ {\omega} ^ {\prime} \circ v _ {\rho (\omega)} \quad \text { and } \quad v _ {\sigma (\omega)} \circ q _ {\omega} = q _ {\omega} ^ {\prime} \circ v _ {\sigma (\omega)}\tag{27}
\]

（换言之，对于所有 \(\lambda \in \mathbf{L}\) , \(v_{\lambda}\) 是关于 \(\mathbf{G}_{\lambda}\) （相应地 \(\mathrm{G}_{\lambda}^{\prime}\) ）上由 \(p_{\xi}\) 和 \(q_{\eta}\) （相应地 \(p_{\xi}^{\prime}\) 和 \(q_{\eta}^{\prime}\) ）（其中 \(\xi\) 和 \(\eta\) 满足 \(\rho(\xi) = \lambda\) 和 \(\sigma(\eta) = \lambda\) ）定义的作用律的一个态射）。则映射

\[
u: (x _ {\lambda}) \mapsto \bigotimes_ {(c, p', q')} v _ {\lambda} (x _ {\lambda})
\]

（从 G 到 \(\bigotimes_{(c, p', q')} \mathrm{G}_{\lambda}'\) ）是 Z-多重线性的，并且满足条件 (26)，因此定义了一个 Z-线性映射（从 \(\bigotimes_{(c, p, q)} \mathrm{G}_{\lambda}\) 到 \(\bigotimes_{(c, p', q')} \mathrm{G}_{\lambda}'\) ），若无混淆之虞，我们将它简记为 \(\bigotimes_{(c)} v_{\lambda}\) 。

我们现在给出如此定义的一般张量积的一个“结合性”性质。设 \((\mathbf{L}_{i})_{1\leq i\leq n}\) 是 L 的一个有限划分；对每个指标 i，令 \(\Omega_{i}\) 表示 \(\Omega\) 中由满足 \(\rho(\omega)\in\mathbf{L}_{i}\) 和 \(\sigma(\omega)\in\mathbf{L}_{i}\) 的元素组成的子集；显然，诸 \(\Omega_{i}\) 两两不相交；我们设 \(\Omega^{\prime}=\Omega-\left(\bigcup_{i}\Omega_{i}\right)\) 。对每个索引 i， \(c^{(i)}\) 将表示映射 \(\omega\mapsto(\rho(\omega),\sigma(\omega))\) （从 \(\Omega_{i}\) 到 \(L_{i}\times L_{i}\) ）；对于 \(\omega\in\Omega_{i}\) ，我们记 \(p_{\omega}^{(i)}\) 和 \(q_{\omega}^{(i)}\) 以代替 \(p_{\omega}\) 和 \(q_{\omega}\) 。于是对每个 i 存在一个“部分”张量积

\[
\mathrm{F} _ {i} = \bigotimes_ {(c ^ {(i)}, p ^ {(i)}, q ^ {(i)})} \mathrm{G} _ {\lambda}.
\]

我们进一步作如下“可置换性”假设：

\begin{enumerate}
\def\labelenumi{(\Alph{enumi})}
\setcounter{enumi}{15}
\tightlist
\item
  若 \(\omega \in \Omega'\) ，则 \(p_{\omega}\) （相应地 \(q_{\omega}\) ）与自同态 \(p_{\xi}\) 和 \(q_{\eta}\) （ \(\mathrm{G}_{\rho(\omega)}\) （相应地 \(\mathrm{G}_{\sigma(\omega)}\) ）的，其中 \(\xi \notin \Omega'\) , \(\eta \notin \Omega'\) 且 \(\rho(\omega) = \rho(\xi) = \sigma(\eta)\) （相应地 \(\sigma(\omega) = \rho(\xi) = \sigma(\eta)\) ））都可交换。
\end{enumerate}

对于每个 \(\omega \in \Omega'\) ，设 \(i\) 是使得 \(\rho(\omega) \in \mathbf{L}_i\) 的索引；然后考虑族 \((v_{\lambda})_{\lambda \in \mathbf{L}_i}\) （其中 \(v_{\rho(\omega)} = p_{\omega}\) ， \(v_{\lambda} = 1_{\mathbf{G}_{\lambda}}\) （对 \(\lambda \neq \rho(\omega)\) ））；假设 (P) 意味着族 \((v_{\lambda})\) 满足条件 (27)（其中 \(p'\) 和 \(p\) 必须换成 \(p^{(i)}\) , \(q'\) 和 \(q\) 换成 \(q^{(i)}\) , \(\omega\) 换成某个元素 \(\xi\) （它遍历 \(\Omega_i\) ））；由此导出一个自同态 \(\bigotimes_{(c^{(i)})} v_{\lambda} = r_{\omega}\) （ \(\mathbf{Z}\) -模 \(\mathbf{F}_i\) 的）。类似地，一个自同态 \(s_{\omega}\) （ \(\mathbf{Z}\) -模 \(\mathbf{F}_j\) 的）由 \(q_{\omega}, j\) 定义，后者是使得 \(\sigma(\omega) \in \mathbf{L}_j\) 的索引；最后令 \(d(\omega) = (i,j)\) 。于是我们可以定义张量积 \(\bigotimes_{(d,r,s)} \mathbf{F}_i\) 以及相应的典范映射

\[
\phi_ {(d, r, s)}: \prod_ {i = 1} ^ {n} \mathrm{F} _ {i} \rightarrow \bigotimes_ {(d, r, s)} \mathrm{F} _ {i}.
\]

另一方面，对每个 i，典范映射

\[
\psi_ {i} = \phi_ {(c ^ {(i)}, p ^ {(i)}, q ^ {(i)})}: \prod_ {\lambda \in L _ {i}} G _ {\lambda} \rightarrow F _ {i};
\]

利用集合乘积的结合性，可导出一个 Z-多重线性映射 \(\psi = \phi_{(d,r,s)}\circ (\psi_i)\) （从 G 到 \(\bigotimes_{(d,r,s)}\mathrm{F}_i\) ）。我们证明有序对 \(\left(\bigotimes_{(d,r,s)}\mathrm{F}_i,\psi\right)\) 是与 \(\left(\bigotimes_{(c,p,q)}\mathrm{G}_{\lambda},\phi_{(c,p,q)}\right)\) 相同的泛映射问题的解，由此将推出存在唯一的 Z-模同构

\[
\theta \colon \bigotimes_ {(c, p, q)} \mathrm{G} _ {\lambda} \rightarrow \bigotimes_ {(d, r, s)} \mathrm{F} _ {i}
\]

，使得 \(\psi = \theta \circ \phi_{(c,p,q)}\) （集合论，IV，§ 3，第 1 号）。

对 \(n\) 作归纳，证明归结为情形 \(n = 2\) ；为简单起见，我们用 \(\mathbf{F}_1 \otimes_{(d)} \mathbf{F}_2\) 和 \(y_1 \otimes_{(d)} y_2\) 分别代替 \(\bigotimes_{(d, r, s)} \mathbf{F}_i\) 和 \(\bigotimes_{(d, r, s)} y_i\) 。考虑从 \(G\) 到 \(\mathbf{F}_1 \otimes_{(d)} \mathbf{F}_2\)

\[
h \colon (x _ {\lambda}) \to \left(\bigotimes_ {(c ^ {(1)})} x _ {\lambda}\right) \underset {(d)} {\otimes} \left(\bigotimes_ {(c ^ {(2)})} x _ {\lambda}\right).
\]

的映射。它显然是 Z-多重线性的；我们证明它对所有 \(\omega\in\Omega\) 满足条件（26）。当 \(\omega\in\Omega_{1}\) 或 \(\omega\in\Omega_{2}\) 时这是显然的；否则，为固定想法起见，假设 \(\rho(\omega)\in\mathbf{L}_{1}\) 和 \(\sigma(\omega)\in\mathbf{L}_{2}\) ，则 \(h\circ\tilde{p}_{\omega}\) 和 \(h\circ\tilde{q}_{\omega}\) 在 \((x_{\lambda})\) 分别为

\[
\left(r _ {\omega} \left(\bigotimes_ {(c ^ {(1)})} x _ {\lambda}\right)\right) \underset {(d)} {\otimes} \left(\bigotimes_ {(c ^ {(2)})} x _ {\lambda}\right) \quad \text { and } \quad \left(\bigotimes_ {(c ^ {(1)})} x _ {\lambda}\right) \underset {(d)} {\otimes} \left(s _ {\omega} \left(\bigotimes_ {(c ^ {(2)})} x _ {\lambda}\right)\right)
\]

处的值，根据 \(\mathbf{F}_1\otimes_{(d)}\mathbf{F}_2\)

的定义也相等。既然如此，设 \(u\) 是一个 \(\mathbf{Z}\) -多重线性映射（从 \(G\) 到一个 \(\mathbf{Z}\) -模 \(H\) ），满足条件（26）；我们将定义一个 \(\mathbf{Z}\) -线性映射 \(v: F_1 \otimes F_2 \to H\) ，使得 \(u = v \circ h\) ，而这将证明我们的断言（重复第 1 号中命题 1 的推论 1 的论证）。对所有 \(z_2 = (x_\lambda)_{\lambda \in L_2}\) ，我们考虑从 \(\prod_{\lambda \in L_1} G_\lambda\) 到 \(H\) 的“部分”线性映射

\[
u \left(\cdot , z _ {2}\right): \left(x _ {\lambda}\right) _ {\lambda \in L _ {1}} \mapsto u \left(\left(x _ {\lambda}\right) _ {\lambda \in L _ {1}}, z _ {2}\right) = u \left(\left(x _ {\lambda}\right) _ {\lambda \in L}\right).\tag{28}
\]

显然它是 Z-多重线性的，并且对 \(\omega\in\Omega_{1}\) 满足条件（26）；因此根据定义，存在一个 Z-线性映射 \(y_{1}\mapsto w_{1}(y_{1},x_{2})\) （从 \(F_{1}\) 到 H 中），使得

\[
w _ {1} \left(\bigotimes_ {(c ^ {(1)})} x _ {\lambda}, z _ {2}\right) = u \left(\left(x _ {\lambda}\right) _ {\lambda \in L _ {1}}, z _ {2}\right)\tag{29}
\]

接下来我们考虑映射

\[
u _ {2}: (x _ {\lambda}) _ {\lambda \in L _ {2}} \mapsto w _ {1} (\cdot , (x _ {\lambda}) _ {\lambda \in L _ {2}})
\]

（一个从 \(\prod_{\lambda \in L_2} G_\lambda\) 到 \(\mathrm{Hom}_{\mathbf{Z}}(\mathrm{F}_1, \mathrm{H})\) 中的映射）；它显然是 \(\mathbf{Z}\) -多重线性的，并且对 \(\omega \in \Omega_2\) 满足条件 (26)，这可由关于 \(u\) 的假设以及关系式 (28) 和 (29)、并注意到形如 \(\bigotimes_{(c^{(1)})} x_\lambda\) 的元素生成 \(\mathbf{Z}\) -模 \(\mathrm{F}_1\) 这一事实得出。因此存在一个 \(\mathbf{Z}\) -线性映射

\[
w _ {2}: \mathrm{F} _ {2} \rightarrow \operatorname{Hom} _ {\mathbf {Z}} (\mathrm{F} _ {1}, \mathrm{H})
\]

，使得

\[
w _ {2} \left(\bigotimes_ {(c ^ {(2)})} x _ {\lambda}\right) = u _ {2} \left(\left(x _ {\lambda}\right) _ {\lambda \in L _ {2}}\right)
\]

或等价地

\[
\left(w _ {2} \left(\bigotimes_ {(c ^ {(2)})} x _ {\lambda}\right)\right) \left(\bigotimes_ {(c ^ {(1)})} x _ {\lambda}\right) = u \left(\left(x _ {\lambda}\right) _ {\lambda \in L}\right).\tag{30}
\]

我们现在考虑，对于 \(y_{1} \in F_{1}, y_{2} \in F_{2}\) ，H的元素

\[
w (y _ {1}, y _ {2}) = (w _ {2} (y _ {2})) (y _ {1}).\tag{31}
\]

显然 \(w\) 是一个 \(\mathbf{Z}\) -双线性映射（从 \(\mathbf{F}_1 \times \mathbf{F}_2\) 到 \(H\) ）。我们进一步证明，对所有 \(\omega \in \Omega'\) ，（为固定思路，设 \(\rho(\omega) \in L_1\) 和 \(\sigma(\omega) \in L_2\) ）

\[
w (r _ {\omega} (y _ {1}), y _ {2}) = w (y _ {1}, s _ {\omega} (y _ {2})).\tag{32}
\]

只需验证此关系当 \(y_{1}\) （相应地 \(y_{2}\) ）具有形式 \(\bigotimes_{(c^{(1)})} x_{\lambda}\) （相应地 \(\bigotimes_{(c^{(2)})} x_{\lambda}\) ）时成立，因为这些元素生成 Z-模 \(F_{1}\) （相应地 \(F_{2}\) ）。但根据定义， \(r_{\omega}\left(\bigotimes_{(c^{(1)})} x_{\lambda}\right) = \bigotimes_{(c^{(1)})} x_{\lambda}'\) ，其中 \(x_{\rho(\omega)}' = p_{\omega}(x_{\rho(\omega)})\) ，而 \(x_{\lambda}' = x_{\lambda}\) （对于 \(\lambda \neq \rho (\omega)\) ， \(\mathbf{L}_1\) 中）；类似地 \(s_{\omega}\left(\bigotimes_{(c^{(2)})}x_{\lambda}\right) = \bigotimes_{(c^{(2)})}x_{\lambda}^{\prime \prime}\) ，其中 \(x_{\sigma (\omega)}^{\prime \prime} = q_{\omega}(x_{\sigma (\omega)})\) ，而 \(x_{\lambda}^{\prime \prime} = x_{\lambda}\) （对于 \(\lambda \neq \sigma (\omega)\) ， \(\mathbf{L}_2\) 中）；利用 (30) 和 (31)，关系式 (32) 随后由 (26) 得出。因此存在一个 Z-线性映射 \(v\) （从 \(\mathrm{F}_1\otimes \mathrm{F}_2\) 到 H 中），使得 \(v(y_{1}\otimes y_{2}) = w(y_{1},y_{2})\) ，并且由 (30) 和 (31) 可知， \(v\circ h = u\) 。

上述定义的一般张量积最重要的特殊情形如下：我们从一族环 \((\mathbf{A}_i)_{1\leqslant i\leqslant n - 1}\) 和一族 \((\mathbf{E}_i)_{1\leqslant i\leqslant n}\) 出发，其中 \(\mathbf{E}_1\) 是一个右 \(\mathbf{A}_1\) -模， \(\mathbf{E}_n\) 是一个左 \(\mathbf{A}_{n - 1}\) -模，并且对于 \(2\leqslant i\leqslant n - 1\) ， \(\mathbf{E}_i\) 是一个 \((\mathbf{A}_{i - 1},\mathbf{A}_i)\) -双模。则上述定义按如下方式应用：L 是集合 \([1,n]\) ， \(\mathrm{G}_i = \mathrm{E}_i\) ， \(\Omega\) 是诸 \(\mathbf{A}_i\) ( \(1\leqslant i\leqslant n - 1\) ) 之和所成的集合。对于 \(\omega \in \mathbf{A}_i\) ( \(1\leqslant i\leqslant n - 1\) )，取 \(\rho (\omega) = i\) ， \(\sigma (\omega) = i + 1\) ， \(p_{\omega}\) 是自同态 \(x\mapsto x\omega\) （ Z-模 \(\mathbf{E}_i\) 的）， \(q_{\omega}\) 是自同态 \(y\mapsto \omega y\) （ Z-模 \(\mathbf{E}_{i + 1}\) 的）；相应的张量积记为

\[
\mathbf {E} _ {1} \otimes_ {\mathbf {A} _ {1}} \mathbf {E} _ {2} \otimes_ {\mathbf {A} _ {2}} \mathbf {E} _ {3} \otimes \dots \otimes_ {\mathbf {A} _ {n - 2}} \mathbf {E} _ {n - 1} \otimes_ {\mathbf {A} _ {n - 1}} \mathbf {E} _ {n}\tag{33}
\]

（一种记号，其中 \(A_{i}\) 有时可省略），该张量积的元素 \(\bigotimes_{(c,p,q)} x_{i}\) ，对于一族 \((x_{i})\) （满足 \(x_{i} \in E_{i}\) ， \(1 \leqslant i \leqslant n\) ），若不会引起混淆，可以写成 \(x_{1} \otimes x_{2} \otimes \cdots \otimes x_{n}\) ；类似的记号也用于 Z-线性映射 \(\bigotimes_{(c)} v_{i}\) 。假设 (P) 对 \([1,n]\) 的每个划分都成立，因为诸 \(E_{i}\) 是双模（ \(2 \leqslant i \leqslant n - 1\) ）。当 n = 3 时，我们由此定义了第 8 号中提到的 Z-模 \(E \otimes_{A} F \otimes_{B} G\) ，并重新得到了第 8 号中的命题 8（第 8 号）。

当每个 \(E_{i}\) 都是多重模（对于 \(2 \leqslant i \leqslant n - 1\) ， \(A_{i-1}\) 是在左边作用的一个环， \(A_{i}\) 是在右边作用的一个环，且对 i = 1 和 i = n 有类似条件），如第 4 号所述，则在 \(E_{1} \otimes_{A_1} E_{2} \otimes \cdots \otimes_{A_{n-1}} E_{n}\) 上定义了关于除 \(A_{i}\) （它们在 \(E_{i}\) ( \(1 \leqslant i \leqslant n\) ) 上作用）之外的所有环的多重模结构。

特别地，设 C 为一个交换环， \((\mathbf{E}_{i})_{1\leqslant i\leqslant n}\) 为一族 C-模。通过给 \(E_{1}\) 和 \(E_{n}\) 两个与给定结构相同的 C-模结构，并给 \(E_{i}\) （对于 \(2\leqslant i\leqslant n-1\) ）三个与给定结构相同的 C-模结构，我们在张量积上定义

\[
\mathbf {E} _ {1} \otimes_ {\mathbf {c}} \mathbf {E} _ {2} \otimes_ {\mathbf {c}} \mathbf {E} _ {3} \otimes \dots \otimes_ {\mathbf {c}} \mathbf {E} _ {n - 1} \otimes_ {\mathbf {c}} \mathbf {E} _ {n}\tag{34}
\]

这样便得到 \(n\) 个彼此相容且实际上相同的 C-模结构，因为对所有 \(\gamma \in \mathbf{C}\) 和 \((x_{i})\in \prod_{i = 1}^{n}\mathbf{E}_{i}\) ，根据定义

\[
\begin{array}{r l} \left(\gamma x _ {1}\right) \otimes x _ {2} \otimes \dots \otimes x _ {n} \\ & = x _ {1} \otimes \left(\gamma x _ {2}\right) \otimes \dots \otimes x _ {n} = \dots = x _ {1} \otimes x _ {2} \otimes \dots \otimes \left(\gamma x _ {n}\right). \end{array}
\]

当我们把张量积（34）称为 C-模时，除非另有说明，我们总是指具有这种结构的张量积；张量积（34）若无混淆，也记作 \(\bigotimes_{1\leqslant i\leqslant n}\mathbf{E}_i\) 。对每个 C-模 G，从 \(\prod_{i=1}^{n}\mathbf{E}_i\) 到 G 中的 Z-多重线性映射，若对每个指标 \(i\) 满足关系

\[
f (x _ {1}, \dots , x _ {i - 1}, \gamma x _ {i}, x _ {i + 1}, \dots , x _ {n}) = \gamma f (x _ {1}, \dots , x _ {n})\tag{35}
\]

对 \(\gamma \in \mathbf{C}\) 和 \((x_{i})\in \prod_{i}\mathbf{E}_{i}\) ，则称为 C-多重线性的，并构成一个 C-模，记为 \(\mathcal{L}_n(\mathbf{E}_1,\dots ,\mathbf{E}_n;\mathbf{G})\) ；张量积（34）的泛性质随后使我们能够定义一个典范的 C-模同构

\[
\mathcal {L} _ {n} \left(\mathrm{E} _ {1}, \dots , \mathrm{E} _ {n}; \mathrm{G}\right)\rightarrow \operatorname{Hom} _ {\mathrm{C}} \left(\mathrm{E} _ {1} \otimes_ {\mathrm{C}} \mathrm{E} _ {2} \otimes \dots \otimes_ {\mathrm{C}} \mathrm{E} _ {n}, \mathrm{G}\right)\tag{36}
\]

它将每个C-多重线性映射f与满足以下条件的C-线性映射g相关联：

\[
f (x _ {1}, \dots , x _ {n}) = g (x _ {1} \otimes x _ {2} \otimes \dots \otimes x _ {n}).
\]

从 \(E_{1} \times \cdots \times E_{n}\) 到 C 中的一个 C-多重线性映射也称为 n-线性形式。

设 \((\mathbf{F}_i)_{1\leqslant i\leqslant n}\) 是另一族 C-模；对于每一个由 \(n\) 个 C-线性映射 \(u_i: \mathbf{E}_i \to \mathbf{F}_i\) 组成的系统， \(u_1 \otimes u_2 \otimes \cdots \otimes u_n\) （也记作 \(\bigotimes_{1\leqslant i\leqslant n} u_i\) ）是一个 C-线性映射，

\[
\mathbf {E} _ {1} \otimes_ {\mathbb {C}} \mathbf {E} _ {2} \otimes \dots \otimes_ {\mathbb {C}} \mathbf {E} _ {n} \quad \text { into } \mathbf {F} _ {1} \otimes_ {\mathbb {C}} \mathbf {F} _ {2} \otimes \dots \otimes_ {\mathbb {C}} \mathbf {F} _ {n}.
\]

进一步， \((u_{1},\ldots,u_{n})\mapsto u_{1}\otimes u_{2}\otimes\cdots\otimes u_{n}\) 是从 \(\prod_{i}\operatorname{Hom}_{\mathbb{C}}(\mathbf{E}_{i},\mathbf{F}_{i})\) 到

\[
\operatorname{Hom} _ {\mathbf {c}} (\mathbf {E} _ {1} \otimes_ {\mathbf {c}} \mathbf {E} _ {2} \otimes \dots \otimes_ {\mathbf {c}} \mathbf {E} _ {n}, \mathbf {F} _ {1} \otimes_ {\mathbf {c}} \mathbf {F} _ {2} \otimes \dots \otimes_ {\mathbf {c}} \mathbf {F} _ {n}).
\]

中的一个 C-多重线性映射。因此，典范地对应于后一个映射的是一个 C-线性映射，称为典范映射

\[
\begin{array}{r l} & \operatorname{Hom} _ {\mathbf {c}} (\mathrm{E} _ {1}, \mathrm{F} _ {1}) \otimes_ {\mathbf {c}} \operatorname{Hom} _ {\mathbf {c}} (\mathrm{E} _ {2}, \mathrm{F} _ {2}) \otimes \dots \otimes_ {\mathbf {c}} \operatorname{Hom} _ {\mathbf {c}} (\mathrm{E} _ {n}, \mathrm{F} _ {n}) \\ & \qquad \to \operatorname{Hom} _ {\mathbf {c}} (\mathrm{E} _ {1} \otimes_ {\mathbf {c}} \mathrm{E} _ {2} \otimes \dots \otimes_ {\mathbf {c}} \mathrm{E} _ {n}, \mathrm{F} _ {1} \otimes_ {\mathbf {c}} \mathrm{F} _ {2} \otimes \dots \otimes_ {\mathbf {c}} \mathrm{F} _ {n}) \end{array}\tag{37}
\]

它推广了第 5 号中对 \(n = 2\) 的情形所定义的映射。

前面看到的一般结合律性质可以在此处特化如下。给定一个划分 \((\mathbf{J}_{k})_{1\leqslant k\leqslant m}\) （ N 的区间 \([1,n]\) 的），对每个 k，令 \(F_{k}\) 为张量积 \(E_{i_{1}}\otimes_{C}E_{i_{2}}\otimes\cdots\otimes_{C}E_{i_{r}}\) ，其中 \((i_{1},\ldots,i_{r})\) 是 \(J_{k}\) 的元素组成的严格递增序列。那么存在一个典范的 C-模同构（称为“结合同构”）

\[
\mathrm{F} _ {1} \otimes_ {\mathbf {c}} \mathrm{F} _ {2} \otimes \dots \otimes_ {\mathbf {c}} \mathrm{F} _ {m} \rightarrow \mathrm{E} _ {1} \otimes_ {\mathbf {c}} \mathrm{E} _ {2} \otimes \dots \otimes_ {\mathbf {c}} \mathrm{E} _ {n}
\]

在上述记号中，它把张量积

\[
y _ {1} \otimes y _ {2} \otimes \dots \otimes y _ {m}, \quad \text { where } y _ {k} = x _ {i _ {1}} \otimes x _ {i _ {2}} \otimes \dots \otimes x _ {i _ {r}},
\]

映射到张量积 \(x_{1} \otimes x_{2} \otimes \cdots \otimes x_{n}\) （其中 \(x_{i} \in E_{i}\) 对所有 i 成立）。

特别是，如果 \(\pi\) 是 \([1, n]\) 的一个排列，写 \(\mathbf{J}_k = \{\pi(k)\}\) （对于 \(1 \leqslant k \leqslant n\) ），我们就得到一个典范的（“交换性”）同构

\[
\mathbf {E} _ {\pi (1)} \otimes_ {\mathbf {C}} \mathbf {E} _ {\pi (2)} \otimes \dots \otimes_ {\mathbf {C}} \mathbf {E} _ {\pi (n)} \rightarrow \mathbf {E} _ {1} \otimes_ {\mathbf {C}} \mathbf {E} _ {2} \otimes \dots \otimes_ {\mathbf {C}} \mathbf {E} _ {n}
\]

它把 \(x_{\pi(1)} \otimes x_{\pi(2)} \otimes \cdots \otimes x_{\pi(n)}\) 映到 \(x_{1} \otimes x_{2} \otimes \cdots \otimes x_{n}\) 。我们常常将这些典范同构之下彼此对应的各种张量积等同起来。

对于 \(1 \leqslant i \leqslant n\) ，进一步假设 \(\mathbf{E}_i\) 有一组基 \((b_{\lambda_1}^{(i)})_{\lambda_i \in \mathbf{L}_i}\) ；通过对 \(n\) 作归纳，由第 7 号命题 7 的推论 2 可知，族 \((b_{\lambda_1}^{(1)} \otimes b_{\lambda_2}^{(2)} \otimes \cdots \otimes b_{\lambda_n}^{(n)})\) （其中 \((\lambda_1, \ldots, \lambda_n)\) 遍历 \(\prod_{1 \leqslant i \leqslant n} \mathbf{L}_i\) ）是 \(\bigotimes_{1 \leqslant i \leqslant n} \mathbf{E}_i\) 的一组基，有时称为所考虑的基 \((b_{\lambda_1}^{(1)})\) 的张量积。

注记. (1) 上述关于交换环上模情形的注记，可以如第 5 号注记 2 那样推广到存在张量积 \(\mathbf{E}_1 \otimes_{\mathbf{A}_1} \mathbf{E}_2 \otimes \cdots \otimes_{\mathbf{A}_{n-1}} \mathbf{E}_n\) 的情形，其中环 \(\mathbf{A}_i\) 不一定是交换的，并且对每个 \(i\) ，存在同态 \(\rho_i: \mathbf{C} \to \mathbf{A}_i\) （同一个交换环 \(\mathbf{C}\) 的），使得：(1) \(\rho_i(\mathbf{C})\) 包含于 \(\mathbf{A}_i\) 的中心；(2) 对于 \(2 \leqslant i \leqslant n - 1\) ， \(\mathbf{E}_i\) 上利用同态 \(\rho_{i-1}\) 和 \(\rho_i\) 得到的 C-模结构是一致的。于是，在 \(\mathbf{E}_1 \otimes_{\mathbf{A}_1} \mathbf{E}_2 \otimes \cdots \otimes_{\mathbf{A}_{n-1}} \mathbf{E}_n\) 上得到一个 C-模结构，并且有类似于 (13)（第 5 号）的典范映射，这些留给读者去描述。

\begin{enumerate}
\def\labelenumi{(\arabic{enumi})}
\setcounter{enumi}{1}
\tightlist
\item
  设 A、B 为两个环，E 为右 A-模， \(E'\) 为左 A-模，F 为右 B-模， \(F'\) 为左 B-模。从 \((\mathbf{E} \otimes_{\mathbf{A}} \mathbf{E}') \times (\mathbf{F} \otimes_{\mathbf{B}} \mathbf{F}')\) 到一个 Z-模 G 中的 Z-双线性映射，与从 \(E \times E' \times F \times F'\) 到 G 中满足条件
\end{enumerate}

\[
\left\{ \begin{array}{l} f (x \lambda , x ^ {\prime}, y, y ^ {\prime}) = f (x, \lambda x ^ {\prime}, y, y ^ {\prime}) \\ f (x, x ^ {\prime}, y \mu , y ^ {\prime}) = f (x, x ^ {\prime}, y, \mu y ^ {\prime}) \end{array} \right.\tag{38}
\]

（对于 \(\lambda\in\mathbf{A}\) , \(\mu\in\mathbf{B}\) , \(x\in\mathbf{E}\) , \(x'\in\mathbf{E}'\) , \(y\in\mathbf{F}\) , \(y'\in\mathbf{F}'\) ）的 Z-多重线性映射 f 成一一对应。本号中给出的一般构造将此证明归结为定义一个典范的 \(\mathbf{Z}\) -模同构（介于 \((\mathbf{E}\otimes_{\mathbf{A}}\mathbf{E}')\otimes_{\mathbf{Z}}(\mathbf{F}\otimes_{\mathbf{B}}\mathbf{F}')\) 与 \(\mathbf{E}\otimes_{\mathbf{A}}\mathbf{E}'\otimes_{\mathbf{Z}}\mathbf{F}\otimes_{\mathbf{B}}\mathbf{F}'\) 之间），这由形如 (33) 的张量积的结合性性质得出。

